% !PRINTED 514
\weberpage{514}
%<<CONT
namely, by the adjunction of algebraic quantities of the simplest
possible nature, to bring about a splitting of the \gesperrt{Galois} resolvent,
and thus a lowering of the degree of the group.

We now put the question thus. In the original
field $\Omega$ the \gesperrt{Galois} resolvent $g(t)$ is irreducible. The
field $\Omega$ is to be extended to another field $\Omega'$ in such a way
that $g(t)$ becomes reducible. Here $\Omega'$ is to be an algebraic field
over $\Omega$, and so by \S.\,143 there must be an algebraic
quantity $\varepsilon$ such that $\Omega' = \Omega(\varepsilon)$. This quantity $\varepsilon$ will be the
root of a certain irreducible equation in $\Omega$, which
we shall denote by
\weq{1}{\chi(\varepsilon) = 0}

Let us assume that in the field $\Omega(\varepsilon)$ the irreducible
factor
\weq{2}{g_1(t) = g_1(t, \varepsilon)}
separates off from $g(t)$, having the root $t = \varrho$. The degree of $g_1(t, \varepsilon)$ we shall
denote by $q$, and the roots by
\weq{3}{\varrho, \varrho_1, \varrho_2 \ldots \varrho_{q-1};}
so that
\weq{4}{g_1(t, \varepsilon) = (t - \varrho)(t - \varrho_1) \ldots (t - \varrho_{q-1}).}
It is now first to be shown that the substitutions
\[ (\varrho, \varrho), (\varrho, \varrho_1), (\varrho, \varrho_2) \ldots (\varrho, \varrho_{q-1}) \]
contained in the group $P$ form a group $Q$. To show this, we set, as
before (\S.\,146),
\[ \varrho_i = \Theta_i(\varrho) \]
and observe that then the equation
\[ g_1[\Theta_i(t), \varepsilon] = 0 \]
has a root, namely $t = \varrho$, in common with $g_1(t)$, and hence,
since $g_1(t)$ is irreducible, is divisible by $g_1(t)$. From this it follows
that, if $\varrho_1$ and $\varrho_2$ are any two of the roots (3), then
$\Theta_i(\varrho_2) = \varrho_3$ is also contained among these roots. But the
substitution $(\varrho, \varrho_3)$ is composed of $(\varrho, \varrho_1)$ and $(\varrho, \varrho_2)$, and
with this the group property of $Q$ is proved. We denote the index of $Q$ by $j$ and set
\weq{5}{p = jq.}
